\documentclass[11pt]{article}
\usepackage{amsmath,amssymb,amsthm,mathtools}
\usepackage[margin=1in]{geometry}
\usepackage{enumitem}

\newcommand{\Pinfty}{P_{\infty}}
\newcommand{\Pinfm}{\mathsf P_{\infty}}
\newcommand{\Mell}{M}
\newcommand{\Uinf}{\mathcal U_{\infty}}
\newcommand{\Cell}{C_{\ell}}
\newcommand{\Chat}{\mathcal C}
\newcommand{\zetaF}{\zeta}

\title{Step 169: Shifted Burnol Packet Factorization Attempt}
\author{Six Birds RH Adequacy Track}
\date{}

\begin{document}
\maketitle

\section{Mode And Target}

This is an attempt-mode computation. The target is the shifted Burnol packet
factorization
\[
  \Mell(B_{\ell,a}u)(s)=\zetaF(s)\alpha_{\ell,a,u}(s),
\]
where \(B_{\ell,a}u\) denotes the inherited shifted packet. No new carrier is
introduced.

\section{T1. Inherited Records Quoted}

\paragraph{R1, Step 165.}
The shifted packet is
\[
  u_{\ell,a}:=B_{\ell,a}u
  :=
  J_a\Pinfty\tau_\ell(I-\Pinfty)u.
\]
The same step records the active multiplier convention
\[
  m_\ell(s)=\exp[-\ell(1/2-s)].
\]

\paragraph{R2, Step 119.}
For a Burnol generator \(g_j\in C_c^\infty([a,A])\), the right Mellin
transform and co-Poisson synthesis are
\[
  \widehat g_j(s)=\int_0^\infty g_j(t)t^{-s}\,dt,
\qquad
  \Chat g_j(t)=\sum_{n\ge1}\frac{g_j(t/n)}{n}-\widehat g_j(1).
\]
With endpoint vanishings imposed, the inherited identity is
\[
  \Mell(\Chat g_j)(s)=\zetaF(s)\widehat g_j(s).
\]

\paragraph{R3, Step 154.}
The shifted commutator form is
\[
  \Cell=(I-\Pinfty)M_{m_\ell}\Pinfty,
\qquad
  r_{\ell,w,k}=\Cell\eta^a_{w,k}.
\]
This records that in the Mellin model \(P_\infty\) is represented by the
projection used with the multiplier \(M_{m_\ell}\).

\paragraph{R4, Step 153.}
The pulled evaluator and transport formulas are
\[
  T_a=\mathcal M_\Gamma J_a\Uinf^{-1},
\qquad
  \Uinf(J_a^*Y^a_{w,k})
  =
  T_a^*\partial_{\bar w}^{\,k}K_a^\Gamma(\cdot,w).
\]
Thus, on the legal transport domain,
\[
  \Mell(J_av)=T_a\Uinf v.
\]

\paragraph{R5, Steps 114--115.}
Step 114 names the sufficient certificate
\[
  \Mell(J_a\Pinfty\tau_\ell(I-\Pinfty)u)(s)
  =
  \zetaF(s)\alpha_{\ell,u}(s).
\]
Step 115 records that a log-shift is a nonvanishing exponential multiplier in
Mellin/Fourier variables and does not itself create a zeta factor.

\section{T2. Derivation}

Let
\[
  F_u:=\Uinf u.
\]

\subsection*{(a) Expand the shifted packet}

By R1 and linearity of \(\Mell\),
\[
\begin{aligned}
  \Mell(B_{\ell,a}u)(s)
  &=
  \Mell(J_a\Pinfty\tau_\ell(I-\Pinfty)u)(s).
\end{aligned}
\]

\subsection*{(b1) Mellin action of \(J_a\)}

R4 defines \(T_a=\mathcal M_\Gamma J_a\Uinf^{-1}\). Therefore for
\[
  v=\Pinfty\tau_\ell(I-\Pinfty)u
\]
in the legal domain,
\[
  \Mell(J_av)=T_a\Uinf v.
\]
Hence
\[
  \Mell(B_{\ell,a}u)
  =
  T_a\Uinf\Pinfty\tau_\ell(I-\Pinfty)u.
\]

\subsection*{(b2) Mellin action of \(P_\infty\)}

The commutator formulas in R3 use the Mellin-side projection \(\Pinfm\).
Thus
\[
  \Uinf\Pinfty\Uinf^{-1}=\Pinfm,
\qquad
  \Uinf(I-\Pinfty)u=(I-\Pinfm)F_u.
\]

\subsection*{(b3) Mellin action of \(\tau_\ell\)}

Using the R1 multiplier convention,
\[
  \Uinf\tau_\ell\Uinf^{-1}=M_{m_\ell},
\qquad
  m_\ell(s)=\exp[-\ell(1/2-s)].
\]
Therefore
\[
  \Uinf\tau_\ell(I-\Pinfty)u
  =
  M_{m_\ell}(I-\Pinfm)F_u.
\]

\subsection*{(c) Compose the actions}

Substitution gives the concrete partial expression
\[
\boxed{
  \Mell(B_{\ell,a}u)(s)
  =
  \left[
  T_a\Pinfm M_{m_\ell}(I-\Pinfm)F_u
  \right](s).
}
\]
This is the full expression obtained from the inherited operator records.

\subsection*{(d) Compare with the co-Poisson identity}

R2 gives a zeta factor only for expressions of the form \(\Mell(\Chat g)\):
\[
  \Mell(\Chat g)(s)=\zetaF(s)\widehat g(s).
\]
Thus the shifted packet would factor if the preceding expression came from a
co-Poisson synthesis:
\[
  J_a\Pinfty\tau_\ell(I-\Pinfty)u=\Chat g_{\ell,a,u}.
\]
In that case
\[
  \Mell(B_{\ell,a}u)(s)
  =
  \zetaF(s)\widehat g_{\ell,a,u}(s),
\qquad
  \alpha_{\ell,a,u}(s)=\widehat g_{\ell,a,u}(s).
\]

\subsection*{(e) Can the correction be absorbed into \(\alpha\)?}

Formally one can write
\[
  \alpha_{\ell,a,u}^{\rm form}(s)
  =
  \frac{\left[T_a\Pinfm M_{m_\ell}(I-\Pinfm)F_u\right](s)}{\zetaF(s)}.
\]
This is a lawful \(\alpha\) only if the numerator vanishes at the nontrivial
zeros to the required multiplicities and satisfies Burnol endpoint, pole, strip,
and carrier conditions. R5 says the shift multiplier is nonvanishing and cannot
itself supply these zeros. Therefore the missing step is not division by
\(\zeta\); it is proving that the transported shifted packet lies in the
co-Poisson range or in the corresponding zeta-divisible Mellin class.

\section{Subclass attempt: u = Cg}

Now specialize to the manager-requested subclass
\[
  u=\Chat g
\]
where \(g\) is a legal Step 119 Burnol generator. Write
\[
  G(s):=\widehat g(s).
\]
By R2,
\[
  \Mell(\Chat g)(s)=\zetaF(s)G(s).
\]
Identifying this with the inherited Mellin realization on the legal subclass,
\[
  F_{\Chat g}:=\Uinf(\Chat g)=M_{\zetaF}G.
\]
Substituting \(F_u=F_{\Chat g}\) into the parent expression gives
\[
\begin{aligned}
  \Mell(B_{\ell,a}\Chat g)
  &=
  T_a\Pinfm M_{m_\ell}(I-\Pinfm)M_{\zetaF}G
  \\
  &=
  T_a\Pinfm M_{m_\ell}
    \bigl(M_{\zetaF}G-\Pinfm M_{\zetaF}G\bigr).
\end{aligned}
\]
The shift multiplier is scalar multiplication, so it commutes with
\(M_{\zetaF}\). Hence
\[
\begin{aligned}
  M_{m_\ell}(I-\Pinfm)M_{\zetaF}G
  &=
  M_{m_\ell}M_{\zetaF}G
  -
  M_{m_\ell}\Pinfm M_{\zetaF}G
  \\
  &=
  M_{\zetaF}(m_\ell G)
  -
  M_{m_\ell}\Pinfm M_{\zetaF}G.
\end{aligned}
\]
After the outgoing projection and the Burnol transport this becomes
\[
\boxed{
\begin{aligned}
  \Mell(B_{\ell,a}\Chat g)
  &=
  T_a\Pinfm M_{\zetaF}(m_\ell G)
  -
  T_a\Pinfm M_{m_\ell}\Pinfm M_{\zetaF}G.
\end{aligned}}
\]
Thus the subclass does add information: the zeta factor enters at the input and
the raw shift preserves it in the first term. The remaining obstruction is not
the shift. It is whether the two projected/transported terms preserve the
zeta-divisible ideal.

\subsection*{Conditional reduction}

The subclass factorization follows from the following more specific identity.

\[
\boxed{\textbf{ZI-COV}^{\rm CP}_{\ell,a}}
\]
There exist legal operators \(A_\infty\) and \(A_{T,a}\) on the relevant
Mellin generator profiles such that
\[
  \Pinfm M_{\zetaF}G=M_{\zetaF}A_\infty G
\]
and, for every profile \(H\) produced below,
\[
  T_a\Pinfm M_{\zetaF}H=M_{\zetaF}A_{T,a}H.
\]
If this zeta-ideal covariance holds, then
\[
\begin{aligned}
  (I-\Pinfm)M_{\zetaF}G
  &=
  M_{\zetaF}(G-A_\infty G),
  \\
  M_{m_\ell}(I-\Pinfm)M_{\zetaF}G
  &=
  M_{\zetaF}\bigl(m_\ell(G-A_\infty G)\bigr),
  \\
  \Mell(B_{\ell,a}\Chat g)
  &=
  M_{\zetaF}
  A_{T,a}\bigl(m_\ell(G-A_\infty G)\bigr).
\end{aligned}
\]
So the conditional quotient is explicit:
\[
  \boxed{
  \alpha_{\ell,a,\Chat g}
  =
  A_{T,a}\bigl(m_\ell(G-A_\infty G)\bigr).
  }
\]
The inherited records do not supply \(A_\infty\), \(A_{T,a}\), or the
zeta-ideal covariance identities. They also do not supply a legal \(g\) for
which a zero-evaluator derivative of either projected term is nonzero. Hence
the subclass result is conditional rather than success or failure.

\section{T3. Upgraded Verdict}

\[
\boxed{\texttt{factorization\_verdict = V\_conditional\_subclass}.}
\]

The parent computation reduced H2 to \(\textbf{CP-RED}_{\ell,a}\). On the
specific subclass \(u=\Chat g\), the calculation sharpens that gap to the
strictly more specific lemma \(\textbf{ZI-COV}^{\rm CP}_{\ell,a}\):
zeta-ideal covariance of \(\Pinfm\) and \(T_a\Pinfm\) on profiles already known
to be \(M_{\zetaF}G\) by Step 119.

\section{T4. Grade}

This upgrade is conditional subclass reduction content. It contains
theorem-grade inherited identities and a new explicit computation on the
co-Poisson subclass. It is not a completed factorization and not a concrete
counterexample.

\section{Inline Nonclaims}

This step does not establish the Riemann Hypothesis. It does not establish
\(\Xi_{\rm BC}=0\) without conditions. It does not prove G2--G5. It does not
prove SL164.1. It does not close any Burnol/Sonine branch. It does not remove
any retained no-go. It does not pivot to a different carrier. It does not claim
that \(\textbf{CP-RED}_{\ell,a}\) is impossible to fill.

\end{document}
